\documentclass{article}
\usepackage{tcee-repaso}
\tcetema[Financiación exterior, deuda externa y ayuda al desarrollo]{3.B.29}{La financiación exterior del desarrollo económico. El problema de la deuda externa. La ayuda al desarrollo.}
\tcefecha{13/11/2026}
\begin{document}
\cabecerarepaso
\begin{columnas}

% ==========================================================================
\bloque{Esquema del cante}
\esquemacante{%
\apartadocante{Introducción: «dale un pez…»}{2,5}
\apartadocante{I. La financiación exterior del desarrollo}{8,5}
\subcante{1. ¿Por qué financiarse fuera?}{3,5}
\subcante{2. El capital no va a los pobres}{1,5}
\subcante{3. Las fuentes}{3,5}
\apartadocante{II. El problema de la deuda externa}{9}
\subcante{1. Sostenibilidad, impago, sobreendeudamiento}{3,5}
\subcante{2. Las crisis y sus respuestas}{3,5}
\subcante{3. Marco institucional y soluciones}{2}
\apartadocante{III. La ayuda al desarrollo}{8,5}
\subcante{1. Concepto y medida}{2,5}
\subcante{2. Evolución y agenda}{2,5}
\subcante{3. ¿Funciona la ayuda?}{3,5}
\apartadocante{Conclusión}{1,5}}
Funciones clave: las 2 brechas; la dinámica de la deuda externa y el saldo que la estabiliza. Gráficos: modelo intertemporal, dinámica de la deuda, curva de \autor{Laffer} de la deuda.

\emph{Pizarra}: introducción y conclusión sin pizarra. I: $I/Y=v\bar\gamma_Y$, la brecha y las 2 brechas con el máximo (en morado, a la izquierda, hasta el final); el intertemporal por capas (autarquía $\to$ apertura $\to$ déficit con $\Delta I$ y $\Delta C$); «58 $\to$ 5» junto al gráfico. II: la dinámica de la deuda y el saldo estabilizador escritos; solo el panel de superávit a mano; la curva de \autor{Laffer} con la quita de $B_1$ a $B_2$; la flecha 1982 $\to$ Brady $\to$ PPME $\to$ Marco Común. III: 2 columnas, «más ayuda» frente a «mejor ayuda», que se llenan con el debate.

% ==========================================================================
\bloque[2,5]{Introducción}
\begin{itemize}
\item \textbf{Enganche}: «dale un pez a un hombre y comerá un día; enséñale a pescar y comerá toda la vida». Sevilla (30/6–3/7/2025), IV Conferencia de Financiación para el Desarrollo: faltan \dato{$\approx$4 billones \$ al año}{ONU}. AOD 2025: \dato{$-23{,}1$\,\% real, 0,26\,\% de la RNB}{OCDE}; servicio de la deuda de renta baja y media \dato{1,4 billones \$}{BM}. ¿Por qué el capital no llega, llega caro o llega regalado y no sirve?
\item \textbf{Relevancia}: sin ahorro exterior la inversión se limita al ahorro propio; \dato{2022-2025: pagaron 436.000 M\$ más de lo que recibieron}{BM}; crisis de los 80 y 90 y ola de deuda tras la pandemia \vertema{3.B.21}; España en el Club de París, los multilaterales y el CAD; TCEE en esos foros \vertema{4.A.24} \vertema{4.A.26}.
\item \textbf{Contextualización}: \concepto{países en desarrollo} (renta baja y media, BM) \vertema{3.B.28}; \concepto{financiación exterior} ($\mathrm{CC}=S-I<0$) \vertema{3.B.11}; \concepto{deuda externa bruta} (pasivos que obligan a pagar; sin acciones ni IED) y neta; \concepto{AOD}. Historia: Plan Marshall y CAD (1961); 3 episodios de deuda (1982; los más pobres hasta 1996-2005; los 2010 con China y la pandemia); 4 conferencias: Monterrey, Doha, Addis Abeba, Sevilla.
\item \textbf{Problemática}: ¿por qué debería ir el capital a los pobres y por qué no va? ¿Cuándo es sostenible la deuda, por qué pagan los Estados y cómo se resuelven las crisis? ¿Qué es la ayuda y funciona?
\item \textbf{Estructura}: \hyperref[rep:I]{I}, todas las fuentes; \hyperref[rep:II]{II}, la que crea deuda; \hyperref[rep:III]{III}, la que no se devuelve.
\end{itemize}

% ==========================================================================
\bloque[8,5]{I. La financiación exterior del desarrollo\label{rep:I}}
\begin{ideaclave}
La teoría dice que el capital debería ir a los pobres y que ganarían los 2; los datos dicen que va poco, y que importa tanto cómo llega como cuánto. Cuanto más estable es y más riesgo asume quien la da, menos vulnerable es el receptor.
\end{ideaclave}
\sub{I.1. ¿Por qué financiarse fuera?\label{rep:I.1}}
\textbf{La brecha de financiación}. Círculo vicioso de la pobreza (\autor{Nurkse}, 1953). \emph{Supuestos}: $v\equiv K/Y$ fija (versión de \autor{Domar} de \autor{Harrod}-\autor{Domar}) \vertema{3.A.43}; $s$ exógena; crecimiento objetivo $\bar\gamma_Y$.
\begin{formula}
\[\frac{I}{Y}=v\,\bar\gamma_Y,\qquad \frac{I-S}{Y}=v\,\bar\gamma_Y-s\]
\end{formula}
\begin{itemize}
\item $v=3$, objetivo 6\,\% $\Rightarrow$ invertir 18\,\%; si ahorra 12\,\%, faltan 6 puntos de PIB al año. Cada punto añade $1/v$ de crecimiento.
\item \emph{2 brechas} (\autor{Chenery}-\autor{Strout}, 1966): ahorro $I-S$ y divisas $M-X$ (bienes de equipo importados); la mayor frena; 3.º límite: capacidad de absorción.
\end{itemize}
\funcionclave{\frac{\mathrm{FE}}{Y}=\max\Big\{v\,\bar\gamma_Y-s,\ \frac{M-X}{Y}\Big\}}
\begin{itemize}
\item Evidencia: no funcionó; Zambia habría superado \dato{20.000 \$}{Easterly, 1999}; tenía $\approx$600. \autor{Domar} renegó del uso. Valoración: rendimientos decrecientes, fungibilidad; no dice si conviene endeudarse.
\end{itemize}
\textbf{El enfoque intertemporal} (\autor{Obstfeld}-\autor{Rogoff}, 1996) \vertema{3.B.12}. Pequeña y abierta, 2 periodos, $U=u(C_t)+\beta u(C_{t+1})$, $F(K)$ cóncava, acceso perfecto a $r^*$.
\begin{itemize}
\item Autarquía en $A$, con $r^A>r^*$ en el país pobre. Con apertura, \concepto{separación} (\autor{Irving Fisher}): $F'(K_{t+1})=r^*$; después reparte el consumo.
\end{itemize}
\begin{formula}
\[\mathrm{CC}_t=S_t-I_t,\qquad \mathrm{XN}_{t+1}=-(1+r^*)\,\mathrm{CC}_t\]
\end{formula}
\graficorepaso[0.65\linewidth]{cc-intertemporal}{Produce en $P$, consume en $C$: déficit en $t$ ($\Delta I+\Delta C$), superávit en $t+1$}
\begin{itemize}
\item El capital debería ir de ricos a pobres y ganan los 2; lo temporal se financia, lo permanente se ajusta. Valoración: sin impago ni restricciones; predice flujos que no se ven.
\end{itemize}
\sub{I.2. El capital no va a los pobres\label{rep:I.2}}
\begin{itemize}
\item \emph{Paradoja de \autor{Lucas}} (1990): India, renta 15 veces menor, capital 0,4 $\Rightarrow$ rendimiento 58 veces mayor; con capital humano, 5; con externalidades, casi nada. Como en el evangelio: «al que tiene, se le dará».
\item Explicaciones: \emph{fundamentales} (PTF, capital humano, infraestructuras, instituciones: \autor{Acemoglu}, \autor{Johnson} y \autor{Robinson}); \emph{fallos del mercado} (impago sin tribunal, información asimétrica de \autor{Gertler}-\autor{Rogoff}, escasez de activos seguros).
\item \concepto{Intolerancia a la deuda} (\autor{Reinhart} et al., 2003); \concepto{puzle de la asignación}: los que más crecen reciben menos (\autor{Gourinchas}-\autor{Jeanne}, 2013); quien más usa ahorro exterior no crece más (\autor{Prasad} et al., 2007).
\end{itemize}
\sub{I.3. Las fuentes\label{rep:I.3}}
\begin{itemize}
\item La pregunta: \emph{quién asume el riesgo}. 2025: remesas \dato{$\approx$700.000 M\$}{BM} > IED \dato{$\approx$494.000 M}{BM}; deuda neta \dato{423.500 M}{BM}; AOD recibida \dato{$\approx$252.000 M (2023)}{BM}.
\item \emph{Sin deuda}: IED ($\geq$10\,\%; estable, comparte riesgo; \autor{Dunning}, OLI; \autor{Helpman}-\autor{Melitz}-\autor{Yeaple}) \vertema{3.B.18}; cartera en acciones (volátil); remesas (estables, anticíclicas, coste >3\,\%); donaciones.
\item \emph{Con deuda}: bancos (petrodólares, 1982); bonos (desde Brady; \concepto{pecado original}: \autor{Eichengreen}-\autor{Hausmann}) \vertema{3.B.17}; oficial bilateral (China) y multilateral \vertema{3.B.22} \vertema{3.B.30}.
\item \emph{Internos e innovadora}: impuestos propios \vertema{4.B.18}; \emph{blended finance}; APP; bonos temáticos; \concepto{canjes de deuda}; clima (\dato{115.900 M\$ en 2022; 300.000 M en 2035}{OCDE; Bakú}) \vertema{3.B.31}.
\end{itemize}
\graficorepaso[0.9\linewidth]{flujos-desarrollo}{Flujos hacia los países de renta baja y media, 2000-2025}

% ==========================================================================
\bloque[9]{II. El problema de la deuda externa\label{rep:II}}
\begin{ideaclave}
Sostenible si el país crece más de lo que le cuesta; un Estado paga por lo que perdería si no pagara; la deuda excesiva perjudica a deudor y acreedor, por eso las crisis acaban en quitas.
\end{ideaclave}
\sub{II.1. Sostenibilidad, impago y sobreendeudamiento\label{rep:II.1}}
\begin{itemize}
\item \concepto{Liquidez}, \concepto{solvencia}, \concepto{sostenibilidad} (\autor{Milesi-Ferretti}-\autor{Razin}, 1996; FMI, 2002). $b$ deuda exterior neta, $d$ déficit comercial (\% PIB):
\funcionclave{b_t=\frac{1+r^*}{1+\gamma_Y}\,b_{t-1}+d_t\ \ \Rightarrow\ \ -d=\frac{r^*-\gamma_Y}{1+\gamma_Y}\,b}
\item Ej.: $b=60$\,\%, $r^*=5$\,\%, $\gamma_Y=2$\,\% $\Rightarrow$ superávit $\approx1{,}8$\,\%. $r^*>\gamma_Y$: \concepto{bola de nieve}; con superávit, $\overline b$ umbral inestable. Concesional: $r^*<\gamma_Y$. En divisas: $\gamma_q$ multiplica.
\end{itemize}
\graficorepaso[0.9\linewidth]{dinamica-deuda-externa}{Con $r^*>\gamma_Y$: déficit, la deuda crece sin límite; superávit, $\overline b$ inestable}
\begin{itemize}
\item \emph{¿Por qué pagan los Estados?} Voluntad, no solo capacidad. Reputación (\autor{Eaton}-\autor{Gersovitz}, 1981): techo de deuda, racionamiento; impago en recesiones (\autor{Arellano}, 2008). Crítica: la reputación sola no basta (\autor{Bulow}-\autor{Rogoff}, 1989) $\Rightarrow$ sanciones directas. Evidencia: costes reales, mayor quita, mayor castigo (\autor{Cruces}-\autor{Trebesch}, 2013).
\item \concepto{Sobreendeudamiento} (\emph{debt overhang}; \autor{Paul Krugman}, 1988; \autor{Jeffrey Sachs}, 1989): la deuda que no se puede pagar es un impuesto sobre la inversión. Quien presta, no cobra; y si cobra, no todo.
\end{itemize}
\graficorepaso[0.6\linewidth]{laffer-deuda}{Curva de \autor{Laffer} de la deuda: a la derecha del máximo, la quita $B_1\to B_2$ sube lo cobrado}
\begin{itemize}
\item Quita buena para los 2; acreedor que resiste (\emph{holdout}). Bolivia 1988: recompra a $\approx$11 centavos, regalo a los acreedores (\autor{Bulow}-\autor{Rogoff}). Más deuda, menos crecimiento (\autor{Pattillo} et al., 2002).
\end{itemize}
\sub{II.2. Las crisis y sus respuestas\label{rep:II.2}}
\begin{itemize}
\item \emph{1982}: petrodólares a tipo variable; \autor{Volcker}; México (8/1982). Bancos de EE.\,UU.: \dato{176\,\% de su capital}{FDIC}. Financiación concertada (liquidez) $\to$ Plan \autor{Baker} (1985, ajuste con crecimiento; Brasil 1987) $\to$ Plan \autor{Brady} (1989, solvencia: bonos con quita y garantía del Tesoro; \dato{17 países hasta 1998}{FMI}; funcionó: \autor{Reinhart}-\autor{Trebesch}, 2016). Década perdida; Consenso de Washington \vertema{3.B.28}.
\item \emph{Los más pobres}: Evian (2003); Jubileo 2000; \concepto{PPME} (HIPC; 1996, reforzada 1999): punto de decisión (deuda >150\,\% de exportaciones en valor actual) y de culminación; \concepto{IADM} (MDRI, 2005): 100\,\% con FMI, AIF y FAfD. \dato{37 de 39 en culminación (Somalia, 12/2023); Sudán en decisión; Eritrea no}{FMI}.
\item \emph{Hoy}: China, mayor acreedor oficial (\autor{Horn} et al., 2021); ISSD (\dato{>12.000 M\$, casi 50 países, 5/2020–12/2021}{G-20}) \vertema{3.B.19}; \concepto{Marco Común} (11/2020: Club de París + China, India, Arabia Saudí; Chad, Zambia, Ghana, Etiopía; lento); Mesa Redonda (2023). Deuda \dato{9,4 billones \$, servicio 13,1\,\% de exportaciones}{BM, 2025}; \dato{9 en sobreendeudamiento, 23 en riesgo alto}{FMI, 3/2026}.
\end{itemize}
\sub{II.3. Marco institucional y soluciones\label{rep:II.3}}
\begin{itemize}
\item Sin tribunal de quiebras. FMI: programa = puerta de entrada \vertema{3.B.22}; \concepto{LIC DSF} con el BM (bajo, moderado, alto, sobreendeudamiento; umbrales por capacidad; revisado en 9/2026, se aplica desde 2027).
\item \emph{Club de París} (1956, Argentina; 22 permanentes, España; Corea y Brasil en 2016): 6 principios (solidaridad, consenso, información, caso por caso, condicionalidad, comparabilidad) y fecha de corte. \emph{Club de Londres}: bancos. \emph{Bonistas}: comités y litigios (Argentina, 2001).
\item Reformas: SDRM (\autor{Anne Krueger}, 2001-2002; abandonado en 2003); \concepto{CAC} (México, 2003); cláusulas climáticas (Sevilla); canjes \vertema{4.A.24}.
\end{itemize}

% ==========================================================================
\bloque[8,5]{III. La ayuda al desarrollo\label{rep:III}}
\begin{ideaclave}
Fuente menor y en retroceso pero decisiva para los más pobres; funciona mejor en programas concretos y evaluados que como relleno de una brecha de ahorro. Importa menos cuánta ayuda se da que cómo y a quién.
\end{ideaclave}
\sub{III.1. Concepto y medida\label{rep:III.1}}
\begin{itemize}
\item Cooperación > ayuda > \concepto{AOD} (CAD): 4 requisitos: oficial; a receptores de la lista del CAD o multilaterales; fin principal el desarrollo; concesional.
\item Desde 2018, \concepto{donación equivalente}: elemento de donación mínimo 45/15/10\,\% (descuento 9/7/6\,\%); antes, 25\,\% con el 10\,\%.
\item Canal (bilateral, multilateral), instrumento, gasto en el donante (refugiados). \concepto{Ayuda ligada}: +15-30\,\% (\autor{Jepma}, 1991); Consenso OCDE (1978), Helsinki (1991), Recomendación del CAD (2001); dilema del prisionero; \dato{81,5\,\% desligada (2022)}{OCDE}. TOSSD.
\end{itemize}
\sub{III.2. Evolución y agenda\label{rep:III.2}}
\begin{itemize}
\item CAD 1961; 0,7\,\% (\autor{Pearson}, 1969; ONU, 1970); fatiga: \dato{0,21\,\% en 2001}{OCDE}; París 2005 (apropiación, alineamiento, armonización, resultados, mutua responsabilidad); Accra, Busan.
\item Financiación del desarrollo: Monterrey 2002, Doha 2008, Addis Abeba 2015, Sevilla 2025 (sin EE.\,UU.; triplicar préstamos multilaterales).
\item Caída: \dato{$-6$\,\% (2024) y $-23{,}1$\,\% (2025): 174.300 M\$, 0,26\,\% de la RNB}{OCDE}; EE.\,UU. $-56{,}9$\,\%; Alemania, 1.\textsuperscript{er} donante; >0,7\,\%: Noruega, Luxemburgo, Suecia, Dinamarca; Reino Unido al 0,3\,\% en 2027. \emph{España}: \dato{4.554 M€, 0,27\,\%}{OCDE, 2025}; Ley 1/2023: 0,7\,\% en 2030 \vertema{4.A.26}.
\end{itemize}
\graficorepaso[0.9\linewidth]{aod-rnb}{AOD del CAD y de España en \% de la RNB, 1960-2025}
\sub{III.3. ¿Funciona la ayuda?\label{rep:III.3}}
\begin{itemize}
\item Motivos del donante: colonia y votos en la ONU (\autor{Alesina}-\autor{Dollar}, 2000). Justificación: brecha o \concepto{gran empujón} (\autor{Rosenstein-Rodan}, 1943; \autor{Sachs}, 2005). Crítica: \autor{Bauer} (1972).
\item \concepto{Paradoja micro-macro}. Por qué no: fungibilidad; enfermedad holandesa (\autor{Rajan}-\autor{Subramanian}, 2011); instituciones (\autor{Deaton}).
\item Econometría: \autor{Burnside}-\autor{Dollar} (2000, buenas políticas $\Rightarrow$ selectividad) $\to$ \autor{Easterly} et al. (2004, desaparece) $\to$ \autor{Rajan}-\autor{Subramanian} (2008, nada robusto) $\to$ \autor{Clemens} et al. (2012, ayuda temprana, efecto modesto) $\to$ \autor{Galiani} et al. (2017, umbral AIF: +0,35 puntos). \emph{Lectura}: positivo pero pequeño.
\item Debate: \autor{Sachs} (más), \autor{Easterly} (buscadores frente a planificadores), \autor{Deaton} y \autor{Moyo} (menos); \autor{Banerjee}, \autor{Duflo} y \autor{Kremer} (Nobel 2019): experimentos aleatorios; límite de validez externa (\autor{Deaton}, 2010) \vertema{4.B.6}.
\end{itemize}

% ==========================================================================
\bloque[1,5]{Conclusión}
\begin{itemize}
\item \textbf{Recapitulación}: el capital va poco a los pobres y llega sobre todo como remesas e IED; deuda sostenible si el crecimiento cubre su coste; quitas que convienen a todos (Brady, PPME, Marco Común); ayuda de efecto pequeño, útil en programas evaluados.
\item \textbf{Relevancia}: ayuda en mínimos desde 2004 y deuda en máximos: la agenda de Sevilla.
\item \textbf{Extensiones}: \vertema{3.B.12} \vertema{3.B.17} \vertema{3.B.22} \vertema{3.B.30} \vertema{3.B.28} \vertema{3.B.31} \vertema{3.B.42} \vertema{4.A.24} \vertema{4.A.26}.
\item \textbf{Opinión}: importa más cómo que cuánto: compartir riesgo, deuda larga y en moneda propia, reestructurar a tiempo, ayuda evaluada.
\item \textbf{Idea final}: el capital no llega porque rinde menos de lo que parece; llega caro como deuda que se come lo pescado; la ayuda como pez no enseña a pescar. La mejor financiación del desarrollo es la que enseña a pescar.
\end{itemize}

% ==========================================================================
\bloque{Las fuentes de un vistazo}
{\small\noindent\begin{tabular}{@{}>{\raggedright\arraybackslash}p{0.22\linewidth}|>{\raggedright\arraybackslash}p{0.1\linewidth}|>{\raggedright\arraybackslash}p{0.27\linewidth}|>{\raggedright\arraybackslash}p{0.2\linewidth}@{}}
\textbf{Fuente} & \textbf{¿Deuda?} & \textbf{Estabilidad} & \textbf{Riesgo}\\ \hline
IED & no & la más estable & inversor\\
Cartera (acciones) & no & volátil; rentas caen en crisis & inversor\\
Remesas & no & estables, anticíclicas & —\\
Donaciones & no & según el presupuesto del donante & donante\\
Bancos y bonos & sí & volátil; parada súbita & país\\
Préstamos oficiales & sí & estables, anticíclicos & país, en condiciones de favor\\
\end{tabular}\par}

% ==========================================================================
\bloque{Indicadores y umbrales}
{\small\noindent\begin{tabular}{@{}>{\raggedright\arraybackslash}p{0.42\linewidth}|c|c|c@{}}
\textbf{LIC DSF (deuda pública externa)} & \textbf{Débil} & \textbf{Media} & \textbf{Fuerte}\\ \hline
VA deuda / PIB (\%) & 30 & 40 & 55\\
VA deuda / exportaciones (\%) & 140 & 180 & 240\\
Servicio / exportaciones (\%) & 10 & 15 & 21\\
Servicio / ingresos fiscales (\%) & 14 & 18 & 23\\
VA deuda pública total / PIB (\%) & 35 & 55 & 70\\
\end{tabular}\par}
PPME: deuda >150\,\% de exportaciones (o 250\,\% de ingresos fiscales en economías muy abiertas). Tasa de renovación <1: la financiación se retira.

% ==========================================================================
\bloque{Las respuestas a las crisis de la deuda}
{\small\noindent\begin{tabular}{@{}>{\raggedright\arraybackslash}p{0.2\linewidth}|>{\raggedright\arraybackslash}p{0.22\linewidth}|>{\raggedright\arraybackslash}p{0.42\linewidth}@{}}
\textbf{Respuesta} & \textbf{Diagnóstico} & \textbf{Qué hizo}\\ \hline
Financiación concertada (1982-85) & liquidez & FMI con ajuste; Club de París aplaza; «dinero nuevo» de los bancos\\
Plan Baker (1985) & liquidez & ajuste con crecimiento; el dinero nuevo no llegó\\
Plan Brady (1989) & solvencia & bonos con quita y garantía del Tesoro; 17 países\\
PPME e IADM (1996-2005) & deuda oficial insostenible & condonación con reformas, también multilateral\\
ISSD y Marco Común (2020) & pandemia; acreedores nuevos & aplazar; reestructurar con China y el Club de París\\
\end{tabular}\par}

% ==========================================================================
\bloque{Concesionalidad de los préstamos (desde 2018)}
{\small\noindent\begin{tabular}{@{}>{\raggedright\arraybackslash}p{0.45\linewidth}|c|c@{}}
\textbf{Receptor} & \textbf{Donación mín.} & \textbf{Descuento}\\ \hline
PMA y otros de renta baja & 45\,\% & 9\,\%\\
Renta media-baja & 15\,\% & 7\,\%\\
Renta media-alta & 10\,\% & 6\,\%\\
Multilaterales & 10\,\% & 5 o 6\,\%\\
\end{tabular}\par}

% ==========================================================================
\bloque{Definiciones}
\begin{itemize}
\item \concepto{Liquidez}: poder refinanciar lo que vence. \concepto{Solvencia}: valor actual de los superávits futuros $\geq$ deuda neta. \concepto{Sostenibilidad}: pagar el servicio sin financiación excepcional ni ajuste brusco (FMI).
\item \concepto{Pecado original}: no poder endeudarse fuera en la propia moneda.
\item \concepto{Intolerancia a la deuda}: crisis con deudas menores en países con historial de impagos.
\item \concepto{Elemento de donación}: nominal menos valor actual de lo que se devuelve, en \% del nominal.
\item \concepto{Ayuda ligada}: obliga a comprar en el país donante.
\item \concepto{Fungibilidad}: la ayuda financia el gasto marginal, no el proyecto.
\item \concepto{Fecha de corte} (Club de París): la deuda posterior no se reestructura.
\item \concepto{CAC}: una mayoría cualificada de bonistas vincula a todos.
\item \concepto{Canje de deuda}: el acreedor perdona a cambio de gasto acordado (desarrollo, naturaleza).
\end{itemize}

% ==========================================================================
\bloque{Datos clave}
\begin{itemize}
\item Deuda externa de renta baja y media: \dato{9,4 billones \$ (+5,5\,\%); servicio 1,4 billones, 13,1\,\% de las exportaciones}{BM, 2025}.
\item Transferencia neta: \dato{$-436.000$ M\$ en 2022-2025}{BM}.
\item LIC DSF: \dato{9 en sobreendeudamiento, 23 en riesgo alto}{FMI, 31/3/2026}.
\item AOD: \dato{174.300 M\$, 0,26\,\% de la RNB, $-23{,}1$\,\%}{OCDE, 2025}; mínimo histórico reciente \dato{0,21\,\% (2001)}{OCDE}.
\item España: \dato{4.554 M€, 0,27\,\% de la RNB}{2025}.
\item Somalia: deuda externa \dato{64\,\% del PIB (2018) $\to$ <6\,\%}{FMI}.
\item Financiación climática: \dato{100.000 M\$ alcanzados en 2022 (115.900)}{OCDE}; NCQG \dato{300.000 M en 2035, llamamiento a 1,3 billones}{CMNUCC}.
\end{itemize}

% ==========================================================================
\bloque{Autores y fechas\label{rep:autores}}
\hito{1943}{\autor{Rosenstein-Rodan}: gran empujón}
\hito{1953}{\autor{Nurkse}: círculo vicioso}
\hito{1956}{Club de París}
\hito{1961}{CAD de la OCDE}
\hito{1966}{\autor{Chenery}-\autor{Strout}: 2 brechas}
\hito{1970}{0,7\,\% (ONU)}
\hito{1981}{\autor{Eaton}-\autor{Gersovitz}}
\hito{1982}{México: crisis de la deuda}
\hito{1988}{\autor{Krugman}: sobreendeudamiento; Bolivia}
\hito{1989}{Plan \autor{Brady}; \autor{Bulow}-\autor{Rogoff}; \autor{Sachs}}
\hito{1990}{\autor{Lucas}: paradoja; Consenso de Washington}
\hito{1996}{PPME (reforzada en 1999)}
\hito{2000}{\autor{Burnside}-\autor{Dollar}}
\hito{2002}{Monterrey; SDRM}
\hito{2005}{IADM; Declaración de París; \autor{Sachs}}
\hito{2015}{Addis Abeba; Agenda 2030}
\hito{2018}{AOD en donación equivalente}
\hito{2020}{ISSD y Marco Común}
\hito{2025}{Sevilla; AOD $-23{,}1$\,\%}

% ==========================================================================
\bloque{Preguntas del tribunal}
\begin{enumerate}
\item \textbf{¿Por qué no va el capital a los pobres?} Rinde menos de lo que parece (fundamentales) y prestarles es arriesgado (impago, información, activos seguros).
\item \textbf{¿Por qué pagan los Estados?} Por las sanciones más que por la reputación (\autor{Bulow}-\autor{Rogoff}).
\item \textbf{¿Cuándo conviene perdonar?} A la derecha del máximo de la curva de \autor{Laffer} de la deuda.
\item \textbf{¿Por qué fracasó Baker y funcionó Brady?} Baker trató un problema de solvencia como de liquidez; Brady redujo la deuda.
\item \textbf{¿Funciona la ayuda?} Efecto positivo pero pequeño en el agregado; sí en programas evaluados.
\item \textbf{¿Qué aporta el Marco Común?} Sienta a China y a los demás acreedores del G-20 con el Club de París; es lento y no llega a la renta media.
\item \textbf{¿Por qué cuesta desligar la ayuda?} Dilema del prisionero entre donantes.
\item \textbf{¿Es malo un déficit exterior?} No en sí: si financia inversión rentable se paga solo; lo temporal se financia, lo permanente se ajusta.
\item \textbf{¿Por qué la deuda sobre el PIB no basta?} No todo el PIB se puede vender fuera: en los pobres, deuda frente a exportaciones e ingresos fiscales.
\item \textbf{¿Qué cambió en 2018 en la AOD?} Los préstamos cuentan solo por su elemento de donación, con umbrales y descuentos según el receptor.
\item \textbf{¿Por qué no hay un SDRM?} Se opusieron EE.\,UU. y los mercados; la vía fue contractual (CAC).
\end{enumerate}

\begin{trampa}
\begin{itemize}
\item 2 brechas: manda la \emph{mayor}; después de los hechos son iguales.
\item Saldo estabilizador: superávit comercial ($-d$) proporcional a $r^*-\gamma_Y$.
\item PPME: 1996, reforzada 1999; IADM: 2005; ISSD: 5/2020–12/2021; Marco Común: 11/2020.
\item Elemento de donación: 45/15/10\,\% (no el 25\,\% anterior a 2018).
\item Club de París: 22 permanentes; España sí, China no. Consenso OCDE: 1978.
\item Plan Baker 1985, Brady 1989; SDRM abandonado (2003); CAC desde 2003 (México).
\end{itemize}
\end{trampa}

\end{columnas}
\end{document}
